\paragraph{Ánh sáng nhìn thấy}
\begin{itemize}
\item Dải bước sóng của ánh sáng nhìn thấy là một phần của thang sóng điện từ.
\item Quang phổ của ánh sáng nhìn thấy là một dải màu biến thiên liên tục từ tím $(\lambda_t=0,38\ \mu m)$ đến đỏ $(\lambda_{\text{đ}}=0,76\ \mu m)$.
\end{itemize}
\paragraph{Các loại bức xạ}
\begin{bang}[tabularx={l|X|X|X}]{CÁC LOẠI TIA}
&\textbf{Tia hồng ngoại}
&\textbf{Tia tử ngoại}
&\textbf{Tia X (Tia Rơn ghen)}
\\\hline
\textbf{Định nghĩa}
&Tia hồng ngoại là các bức xạ không nhìn thấy có bước sóng \textbf{dài hơn} \bth{0,76\microm} đến vài khoảng milimet .
&Tia tử ngoại là các bức xạ không nhìn thấy có bước sóng \textbf{ngắn hơn} \bth{0,38\microm} đến cỡ $10^{-8}\ m$ .
&Bức xạ có bước sóng từ $10^{-11}\ m$ đến $10^{-8}\ m$.
\\\hline
\textbf{Nguồn phát}
&Mọi vật có nhiệt độ cao hơn 0 K đều phát ra tia hồng ngoại.
&
Các vật có nhiệt độ trên 2000\degree C: đèn hơi thủy ngân, hồ quang điện, mặt trời,\ldots\newline
%$\bullet$ Trong chùm ánh sáng Mặt Trời có 9\% năng lượng thuộc về tia tử ngoại.
&Dùng ống Cu-lit-giơ.
\\\hline
\textbf{Đặc điểm}
&
$\bullet$ Có tác dụng nhiệt mạnh.\newline
$\bullet$ Gây ra một số phản ứng hóa học, tác dụng lên một số loại phim ảnh.\newline
$\bullet$ Có thể biến điệu được như sóng điện từ cao tần.\newline
&
$\bullet$ Tác dụng mạnh lên phim ảnh, làm ion hóa không khí và nhiều chất khí khác.\newline
$\bullet$ Kích thích sự phát quang, gây ra một số phản ứng hóa học.\newline
$\bullet$ Bị thủy tinh, nước hấp thụ rất mạnh.\newline
$\bullet$ Hủy diệt tế bào da, làm hại mắt, diệt khuẩn, diệt nấm mốc,\ldots\newline
&
$\bullet$ Có khả năng đâm xuyên mạnh.\newline
$\bullet$ Có tác dụng mạnh lên phim ảnh, làm ion hóa không khí.\newline
$\bullet$ Có tác dụng làm phát quang một số chất.\newline
$\bullet$ Tia X có tác dụng sinh lí mạnh, hủy diệt tế bào, diệt vi khuẩn,\ldots
\\\hline

\textbf{Ứng dụng}
&
$\bullet$ Sưởi ấm, sấy khô quần áo,\ldots \newline
$\bullet$ Điều khiển từ xa.\newline
$\bullet$ Chụp ảnh bề mặt của Trái đất qua vệ tinh.\newline
$\bullet$ Ứng dụng đa dạng trong lĩnh vực quân sự.
&Khử trùng nước, thực phẩm và dụng cụ y tế, chữa bệnh (bệnh còi xương), tìm vết nứt trên bề mặt kim loại,\ldots
&
$\bullet$ Chiếu điện, chụp điện.\newline
$\bullet$ Kiểm tra chất lượng các vật đúc, tìm các vết nứt, các bọt khí bên trong các vật bằng kim loại.
\end{bang}
\paragraph{Thang sóng điện từ}
\begin{center}
includegraphics[width=8cm]{Hinh/songdientu}
\end{center}
\Binhluan{
Sóng vô tuyến, tia hồng ngoại, ánh sáng nhìn thấy, tia tử ngoại, tia Rơn-ghen, tia gam-ma có cùng bản chất là sóng điện từ. Chúng tạo thành dải bước sóng giảm dần (tần số tăng dần)}
{Các loại sóng điện từ \textit{không có một ranh giới nào rõ rệt}}

